\cotefile{156-7}
\scanpage{42}\mspage{41}
\deleted{NB Cet axiome \unsure{ne|en} fait \unsure{intervenir|intervient} $\mathrel{\rlap{$>$}<}$ que par l'intermédiaire de la relation $\parallel$, et \unsure{$=$|même} seulement de la relation $X'\parallel Y'$ entre éléments de $\Lambda$.}
% slanted
\margin{\uline{$M\Lambda4$} $X\mathrel{\rlap{$>$}<} Y$ et $X'\leq X,\ Y'\leq Y$ $\Longrightarrow X'\mathrel{\rlap{$>$}<} Y'$}

\uline{Lemme} $A'$ comme lemme $A$, en supposant seulement $X\mathrel{\rlap{$>$}<} Y$.

Mais on a besoin dans le raisonnement (p. 30) de la relation
\[
\mathrm{omb}_{\Lambda}(X)\cap \mathrm{omb}_{\Lambda}(Y)=\mathrm{omb}_{\Lambda}(L)\overset{\mathrm{def}}{=}\bigcup_{Z\in\widetilde X\cap\widetilde Y}\mathrm{omb}_{\Lambda}Z
\]
ce qu'il faut \unsure{poser|pour} ici \unsure{en|au} \unsure{axiome|minimum}

\[
\left|\begin{aligned}
&\text{\uline{$M\Lambda$}\deleted{\unsure{5}}$\unsure{4''|4'}$}\qquad\text{Si }X\mathrel{\rlap{$>$}<} Y,\quad\text{et si }Z\ll X,\ Z\ll Y,\ \text{alors}\\
&\exists\,Z'\in\widetilde X\cap\widetilde Y\quad\text{avec }Z\ll Z',\quad\text{i.e.}\quad Z\ll Z'\leq X,Y .
\end{aligned}\right.
\]
% slanted
\margin{On peut se borner au cas $Z\in\Lambda$}

(Propriété \unsure{restrictive} sur les \deleted{\unsure{couples}}\inserted{paires} $(X,Y)$ avec $X\mathrel{\rlap{$>$}<} Y$)

% slanted
\margin{\unsure{lemme} 6 p. 32}
\uline{Lemme} $\unsure{B}$ Soient $Z\overset{\circ}{\ll} X,\ Z\overset{\circ}{\ll} Y$, avec $X\mathrel{\rlap{$>$}<} Y$. Alors $X=Y$ (i.e. on a $A\!t_{M}\,4$ b))

\deleted{\unsure{Cor}}\inserted{\uline{Cor}} Plus gén, si \deleted{$Z\overset{\circ}{\ll} X,\ Z\ll Y,$} $X\mathrel{\rlap{$>$}<} Y$ \unsure{et} $Z\overset{\circ}{\ll} X$, \deleted{alors} $Z\ll Y$, alors $X\leq Y$.

\deleted{Utilisant $M\Lambda\,2$ (i.e. $\exists\,Z_0\in\Lambda,\ Z_0\overset{\circ}{\ll} Z$ donc $Z_0\overset{\circ}{\ll} X$ et $Z_0\ll Y$), \unsure{pour conclure donc} \unsure{la condition des}}

\deleted{\uline{Lemme} $C'$ Soient $X\unsure{,}\unsure{Y|X'}$ \inserted{\unsure{$(\ll X)$}} avec $X\mathrel{\rlap{$>$}<} \unsure{Y'|X'}$. \inserted{\unsure{Cond. équiv}} \deleted{\unsure{et} $Z\ll X$}. \deleted{\unsure{Conditions équivalentes}} \unsure{(i)} $Y\ll X'$ $X$ \unsure{(bis)}.}

En effet, soit $Z_0\in\Lambda,\ Z_0\overset{\circ}{\ll} Z$, alors $Z_0\overset{\circ}{\ll} Z\overset{\circ}{\ll} X$ donc $Z_0\overset{\circ}{\ll} X$, \deleted{\unsure{et}} i.e. $Z_0\in \mathrm{omb}_{\Lambda}(X)^\circ$ $Z_0\ll Z\ll Y$ donc $Z_0\ll Y$, i.e. $Z_0\in \mathrm{omb}_{\Lambda}(Y)$, donc $\mathrm{omb}_{\Lambda}(X)^\circ\cap \mathrm{omb}_{\Lambda}(Y)\ne\varnothing$, donc $X\leq Y$ par lemme $\Lambda'$. OK.
\scanpage{43}\mspage{42}
\uline{Proposition} Propriétés caractéristiques d'un système $(M,\leq,\ll,\mathrel{\rlap{$>$}<},\Lambda\subset M)$, pour provenir d'un atlas avec $\Lambda$ dense dans $M$.

% slanted
\margin{Axiome trivial}
\[
\deleted{\unsure{A}}M\underset{\Lambda}{\Lambda}\,0\qquad X\leq Y\quad\Longrightarrow\quad X\ll Y
\]

% slanted
\margin{Axiome du raffinement \unsure{induit|indirect} \deleted{(Mais sera contenu dans $M\Lambda\,4$ et $M\Lambda\,5$).}}
\[
\begin{aligned}
M\underset{\Lambda}{\Lambda}\,1\qquad&\text{Si}\qquad Y\ll \begin{matrix}X\\[-.4ex]\unsure{\rotatebox{90}{$\leq$}|\rotatebox{90}{$\ll$}}\\[-.4ex]X'\end{matrix},\quad \text{\uline{$Z\in\Lambda$}}\quad\text{et}\quad Z\ll Y,\ Z\ll X',\\
&\text{alors}\quad \exists\,\unsure{\mathord{.}}\,Y'\,\deleted{\unsure{\ll}}\ \text{avec}\quad Y'\leq Y,\ Y'\ll X',\ Z\ll Y'\mathpunct{:}
\end{aligned}
\]
\[
\begin{aligned}
\mathrm{omb}_{\Lambda}(Y)\cap \mathrm{omb}_{\Lambda}(X')&=\bigcup_{\substack{\deleted{\unsure{\varnothing}}\\Y'\in\deleted{\unsure{\widetilde Y\cap\widetilde X'}}\inserted{\widetilde Y\cap\mathrm{Omb}(X')}}}\mathrm{omb}_{\Lambda}(Y')\\
\deleted{\unsure{\text{et}}}\qquad&\overset{\text{si }Y\leq X}{=}\bigcup_{Y'\in\widetilde Y\cap\widetilde{X'}}\mathrm{omb}_{\Lambda}(Y')
\end{aligned}
\]
\margin{grande ombre}

% slanted
\margin{Axiome de densité de $\Lambda$}
\[
M\Lambda\,2\qquad \forall\,\unsure{x|X}\in M,\quad\exists\,Z\in\Lambda\quad\text{avec}\quad Z\overset{\circ}{\ll} X
\]

% slanted
\margin{Axiome de la \unsure{stricte voisinance|stricte minimalité}}
\[
\begin{aligned}
M\Lambda\,3\qquad&\forall\,X\in M,\ Z\in\Lambda,\ \text{avec}\ Z\ll X,\ \exists\,X'\leq X\ \text{avec}\\
&Z\overset{\circ}{\ll} X'
\end{aligned}
\]

\[
\begin{aligned}
M\Lambda\,4\qquad &\inserted{\mathrel{\rlap{$>$}<}\ \text{relation réflexive symétrique, et}}\\
&X\mathrel{\rlap{$>$}<} Y,\quad X'\leq X,\quad Y'\leq Y\quad\Longrightarrow\quad X'\mathrel{\rlap{$>$}<} Y'
\end{aligned}
\]

\[
\begin{aligned}
M\Lambda\,\deleted{\unsure{4}}5\qquad &\text{Si}\quad X\mathrel{\rlap{$>$}<} Y\quad\text{et}\quad Z\in\Lambda\quad\text{tels que}\quad Z\ll X,\ Z\ll Y,\\
&\text{alors}\quad\exists\, Z'\in M\quad\text{avec}\quad Z\ll Z'\leq X,Y\\
&\text{i.e.}\quad \mathrm{omb}_{\Lambda}(X)\cap \mathrm{omb}_{\Lambda}(Y)=\bigcup_{Z'\in\widetilde X\cap\widetilde Y}\mathrm{omb}_{\Lambda}(Z')
\end{aligned}
\]
% slanted
\margin{Cette \unsure{condition} implique la condition \unsure{en apparence|apparemment} plus forte, où on ne \unsure{supp.} plus $Z\in\Lambda$.}

C'est ces \inserted{deux} derniers axiomes seulement qui font intervenir $X\mathrel{\rlap{$>$}<} Y$. \deleted{\unsure{$M$}}L'axiome $M\Lambda\,4$ est un axiome de saturation. Pour $M\Lambda\,5$, \deleted{soit} Notons $\unsure{\mathcal R|\mathscr R}$ la relation (définie en termes de $(M,\leq,\ll)$ seulement, sans \unsure{réf} à $\Lambda$ ni à $\mathrel{\rlap{$>$}<}$) :
\[
\begin{aligned}
\mathcal R(X,Y)\quad\Longleftrightarrow\quad &\forall\, Z\in M\ \text{avec}\ Z\ll X,\ Z\ll Y,\\
&\exists\,Z'\in M\ \text{avec}\ Z\ll Z'\leq X,Y\\
&\text{i.e.}\quad\mathrm{Omb}(X)\cap\mathrm{Omb}(Y)=\bigcup_{Z'\in\widetilde X\cap\widetilde Y}\mathrm{Omb}(Z')
\end{aligned}
\]
\scanpage{44}\mspage{43}
\inserted{Cette relation est réflexive symétrique.} On vérifie \deleted{tout} de suite (en utilisant $M\Lambda\,1$ dans le cas de $Y\leq\begin{matrix}X\\[-.4ex]\rotatebox{90}{$\leq$}\\[-.4ex]X'\end{matrix}$) que \deleted{cette \unsure{relat}} \deleted{$\unsure{\mathcal R(X|\mathscr R(X}$} si $X'\leq X,\ Y'\leq Y$, \deleted{alors} $\mathcal R(X,Y)\Longrightarrow\mathcal R(X',Y')$ \deleted{\unsure{et l'}} \unsure{le'|l'}axiome $\unsure{ML|M\Lambda}\,5$ est \deleted{\unsure{de}} \unsure{satisfait|toujours} \unsure{tautologiquement|tautologique} pour cette condition. Ainsi l'axiome $M\Lambda\,5$ prend la forme
\[
X\mathrel{\rlap{$>$}<} Y\quad\Longrightarrow\quad\mathcal R(X,Y)
\]
et on voit donc que les relations cherchées $\mathrel{\rlap{$>$}<}$ sont les relations \inserted{réfl. sym.} ``saturées pour $\leq$''\unsure{)|,} qui impliquent $\mathcal R(X,Y)$. La relation admissible $\mathrel{\rlap{$>$}<}$ la plus large sera obtenue, en prenant $\mathcal R$ elle-\unsure{$=$|même}.

Dans le cas \inserted{d'un atlas} ensembliste spacieux, on a $\mathrel{\rlap{$>$}<}=\mathcal R$, mais pas dans d'autres cas intéressants. Ainsi on note que $\mathcal R(x,y)$ tjrs vrai si $x,y$ sont deux \unsure{$L$} --- or dans les cas non ensemblistes, il \unsure{pourra|arrivera} exister des lieux $x,y$ qui ne soient compatibles.

Je voudrais maintenant regarder le cas où $\mathrel{\rlap{$>$}<}$ peut se définir par la connaissance de $\parallel$ sur $\Lambda$, par une condition

\uline{$At\,C$}$_{\Lambda}$ (\unsure{espac|spac|espace}\deleted{\unsure{e}}) Soient $X,Y\in M$, \deleted{Alors $X\mathrel{\rlap{$>$}<} Y$}\inserted{$L=\widetilde X\cap\widetilde Y$.}
\[
\left|\begin{aligned}
X\mathrel{\rlap{$>$}<} Y\Longleftrightarrow{}&\deleted{\forall\,X',Y'\in\Lambda,\ X'\ll X,\ Y'\ll Y}\\
&\deleted{\forall\,X'\in \mathrm{omb}_{\Lambda}(X)\smallsetminus \mathrm{omb}_{\Lambda}(L)}\\
&\deleted{Y'\in \mathrm{omb}_{\Lambda}(Y)\smallsetminus \mathrm{omb}_{\Lambda}(L)}\\
&\inserted{\forall\,X'\in \mathrm{omb}_{\Lambda}X,\ Y'\in\mathrm{Omb}_{\Lambda}Y\ \text{avec }X'_L=Y'_L=\varnothing,}\\
&\text{on a}\quad X'\parallel Y'
\end{aligned}\right.
\]
(où on a posé, \deleted{pour $L\unsure{\subset|\in}M,\ \mathrm{omb}_{\Lambda}(L)=\bigcup_{Z\in L}\mathrm{omb}_{\Lambda}(Z)$} \inserted{$L=\widetilde X\cap\widetilde Y$, $X'_L=\deleted{\unsure{\widetilde X}}\widetilde{X'}\cap\mathrm{Omb}_{\deleted{\unsure{\Lambda}}}(L)=\left\{Z\,\middle|\,\begin{aligned}&Z\leq X'\\&\exists\,S\in L\\&\deleted{\unsure{\text{avec}}}Z\ll S\end{aligned}\right\}$}\unsure{\}}.

\uline{Cor} Soient $X,Y\in M$. \deleted{tels que} \deleted{Alors $X'$} \inserted{Alors} $X\parallel Y$ \deleted{i.e. $X\deleted{\unsure{\parallel}}\mathrel{\rlap{$>$}<} Y$ et $L=\varnothing$}. \deleted{Si $X'$} $\forall\, X'\in \mathrm{omb}_{\Lambda}(X)\quad\forall\,Y'\in \mathrm{omb}_{\Lambda}Y$ on a $X'\parallel Y'$.
\scanpage{45}\mspage{44}
% slanted
\margin{On veut aussi \uline{$At\,D$}}
% Long diagonal cancellation strokes traverse the main text.
\deleted{La question, c'est quand un système $(M,\leq,\ll,\Lambda\subset M,\underset{\Lambda}{\parallel})$ provient d'une structure d'\unsure{Atelier|Atelies} \inserted{avec $\Lambda\subset M$ dense,} satisfaisant $At\,C_{\Lambda}$ (\unsure{spac})\unsure{.} \inserted{$At\,D$} \deleted{\unsure{Cette}} Les axiomes $M\Lambda\,0$ \unsure{:|à} $M\Lambda\,3$ \deleted{restent} ne \unsure{varieront|varient}. \deleted{Pour} On doit avoir \unsure{$\checkmark$}}

\deleted{\uline{$M\Lambda'$}4 \inserted{La relation} $\parallel$ \inserted{sur $\Lambda$} est symétrique et antiréflexive ($\deleted{X}Z\parallel\deleted{Y}\unsure{Z'}$ implique $\deleted{X}Z\ne\deleted{Y}Z'$). \deleted{\unsure{Si on a} $Z\parallel Z'$} Si $X\parallel Y$, et $X'\ll X,\ Y'\ll Y$, alors $X'\parallel Y'$.}

\deleted{\uline{$M\Lambda'$}5 Soient \deleted{$\unsure{Y\leq X},\quad X',Y'\in\Lambda,\ \text{tels que}$ \unsure{soit} $L=\widetilde X\cap\widetilde Y$, \unsure{$X$ et} $X',Y'\in\Lambda$ avec $X'\ll X,\ Y'\ll Y$} \deleted{\inserted{$X,Y\in M$, tels que il existe $Z$ avec $X,Y\leq Z$.}} $Y\leq X$, $\unsure{\mathord{.}}Y'\in \mathrm{omb}_{\Lambda}(Y),\ X'\in \mathrm{omb}_{\Lambda}(X)$. Supposons $X'_Y=\varnothing$ [i.e. $\mathrm{omb}_{\Lambda}(Y)\cap \mathrm{omb}_{\Lambda}(X')=\varnothing$, ou encore \deleted{$\widetilde{X'}\cap\unsure{\mathrm{omb}}$} $\mathrm{Omb}_{\unsure{\Lambda}}(Y)\cap\widetilde{X'}=\varnothing$ ou encore $\mathrm{Omb}(Y)\cap\mathrm{Omb}(X')=\varnothing$ --- \deleted{\unsure{tout}} ça revient au \unsure{$=$|même}]. Alors $X'\parallel Y'$.}
% slanted
\margin{\uline{NB} On définit \deleted{\unsure{$X'_Y=$}} $X'_Y=\widetilde{X'}\cap\mathrm{Omb}(Y)$}

\deleted{Vérifions d'abord les équivalences entre $[\,]$ :}

\deleted{\uline{Lemme} (\deleted{\unsure{uniquement}} indépendamment de la donnée de $\parallel$ ou de $\mathrel{\rlap{$>$}<}$ --- ne dépend que de $\leq,\ \ll,\ \Lambda$ satisfaisant $M\Lambda\,0$ : $M\Lambda\,3$) Soient $Y\leq\begin{matrix}X\\[-.4ex]\rotatebox{90}{$\ll$}\\[-.4ex]X'\end{matrix}$. Conditions \inserted{$i,i'\ \unsure{(ii')|(ii'')}$} équivalentes \inserted{et impliquent \unsure{$(i0)$|$(ii)$}} :}
\[
\begin{aligned}
&\deleted{\unsure{(i)}\quad \mathrm{omb}_{\Lambda}(Y)\cap \mathrm{omb}_{\Lambda}(X')=\varnothing}\\
&\deleted{\bigl[\unsure{(ii)|(i0)}\quad \mathrm{omb}_{\Lambda}(Y)\cap\widetilde{X'}=\varnothing\bigr]}\\
&\deleted{\unsure{(i')|(i'')}\quad\mathrm{Omb}_{\unsure{M|\Lambda}}(Y)\cap\mathrm{Omb}(X')=\varnothing}\\
&\deleted{\unsure{(ii')|(i'')}\quad\underbrace{\mathrm{Omb}(Y)\cap\widetilde{X'}}_{X'_Y}=\varnothing}
\end{aligned}
\]
\margin{\deleted{$\leftarrow$ est impliqué par les autres, et n'est équivalent que si on suppose $\forall\,X\in M,\ \exists\,Z\in\Lambda$ tel que $Z\leq X$ ($\Lambda$ ``\uline{très} dense'' dans $M$)}}

\deleted{\uline{\unsure{Dém}} \deleted{$(i)\Longrightarrow(ii)$ \unsure{tautologiquemt|tautologique}}, \unsure{$(ii')$}$\Longrightarrow(i')$ : \deleted{\unsure{en}} Supposons \unsure{$(ii')$}, prouvons $(i)$, par $M\Lambda\,1$ on a $\mathrm{omb}_{\Lambda}(Y)\cap \mathrm{omb}_{\Lambda}(X')$}
\[
\deleted{{}=\bigcup_{Y'\in\widetilde{X'}\cap}}
\]
% The lower scan edge cuts the union off after the intersection sign.
\scanpage{46}\mspage{4\deleted{4}5}
\uline{24 Juin}

Soit $\Lambda\subset M$, \unsure{fermé|fermée} pour $\leq$. On introduit des variantes des conditions \unsure{envisagées|analogues} pour $L$ :

% slanted
\margin{densité de $\Lambda$ $\longrightarrow$}
\[
\text{\uline{$At\,\Lambda\,1$}}\qquad \forall\,X\in M,\quad \mathrm{omb}_{\Lambda}(X)^\circ\ne\varnothing\quad\text{i.e.}\quad \exists\,Z\in\Lambda,\ Z\overset{\circ}{\ll} X
\]

% slanted
\margin{Si on a $At\,\Lambda\,1$ et $At\,\Lambda\,2$, on dit que $\Lambda$ est très dense}
\[
\begin{aligned}
\text{\uuline{$At\,\Lambda\,2$}}\qquad \forall\,X,Y\in M,\quad X\parallel Y\quad\Longleftrightarrow\quad &\forall\,X'\in \mathrm{omb}_{\Lambda}(X),\ Y'\in \mathrm{omb}_{\Lambda}(Y)\\
&X'\parallel Y'
\end{aligned}
\]

% slanted
\margin{Si $\Lambda=L$, c'est \unsure{réduit|ident.|évidemt} à $At\,L\,3$}
\uuline{$At\,\Lambda\,3$} \uline{ou $At\,D\Lambda$} Soient $X,Y\in M$ \deleted{tels que il existe} \inserted{avec} $\deleted{\unsure{Z}}X\mathrel{\rlap{$>$}<} Y$, \deleted{\unsure{$X'$} \unsure{et}} $L=X\cap Y,\ X'\in \mathrm{omb}_{\Lambda}(X),\ Y'\in \mathrm{omb}_{\Lambda}(Y)$. Alors
\[
X'_L\mathrel{\underset{\deleted{\unsure{\Lambda}}}{\parallel}}Y'_L\quad\Longrightarrow\quad X'\parallel Y'.
\]

\uline{$At\,C\Lambda$} (spac) Soient $X,Y\in M$, \deleted{Supposons que} \inserted{$L=X\cap Y\deleted{\overset{\mathrm{def}}{=}\unsure{\widetilde X}}$} définie par $\widetilde L=\widetilde X\cap\widetilde Y$). Alors
\[
X\mathrel{\rlap{$>$}<} Y\quad\Longleftrightarrow\quad
\forall\begin{aligned}&X'\in \mathrm{omb}_{\Lambda}(X)\\&Y'\in \mathrm{omb}_{\Lambda}(Y)\end{aligned}
\qquad X'_L=Y'_L=\varnothing\quad\Longrightarrow\quad X'\parallel Y'
\]

\uline{$At\,C\Lambda$} (pol) Soient $X,Y\in M,\ L=X\cap Y$. Alors
\[
X\mathrel{\rlap{$>$}<} Y\Longleftrightarrow\left\{\begin{aligned}
&(i)\quad\text{Comme dans }At\,C\Lambda\ \text{(spac) ci-dessus}\\
&(ii)\quad L=\varnothing\ \text{ou } L\in M
\end{aligned}\right.
\]

On a encore les implications (comparer p. 22)
\diagram{\begin{tikzcd}[ampersand replacement=\&,column sep=1.3em,row sep=1.1em]
{} \& \Lambda3+C\Lambda\mathord{*}\arrow[r,Rightarrow] \& C\deleted{\unsure{*|\Lambda}}\arrow[d,Rightarrow] \& {}\\
{} \& \Lambda2+\Lambda3\arrow[r,Rightarrow]\arrow[d,Rightarrow] \& D\arrow[d,Rightarrow]\arrow[r,Rightarrow] \& \Lambda3\\
D_0+\Lambda1+C\Lambda\mathord{*}\arrow[r,Rightarrow] \& \Lambda2\arrow[r,Rightarrow] \& D_0 \& {}
\end{tikzcd}}
\[
\Lambda2+C\Lambda\mathord{*}\quad\Longleftrightarrow\quad\Lambda2+C\mathord{*}
\]
où on a laissé tomber \unsure{le|les} sigle ``$At$'' dans la désignation des axiomes, et où $C\Lambda\mathord{*}$ désigne indifféremment $At\,C\Lambda$ (\unsure{spac}) ou $At\,C\Lambda$ (pol), et \deleted{et $\unsure{CBS|CDS}$ désigne $At\,C$ (\unsure{le|la} \unsure{B|S} suggère : ``\unsure{sous-figures|bouffe-figures}'' \inserted{\unsure{il bouffe des|il bouff. des}}).} $C\mathord{*}$ est $At\,C$ (\unsure{spac}) ou $At\,C$ (pol) (mais toujours le \unsure{$=$|même})
\scanpage{47}\mspage{46}
Un ensemble \uline{complet} de conditions \unsure{de} compatibilité-\unsure{disjonction|des parties} impliquant $\Lambda$ \deleted{\unsure{peut}} (impliquant donc tous les axiomes pertinents
\diagram{\begin{tikzcd}[ampersand replacement=\&,column sep=.22em,row sep=.55em]
C, \& D,\arrow[d,Rightarrow] \& C\mathord{*},\ \Lambda1,\ \Lambda2,\ \Lambda3,\ C\Lambda\mathord{*}\\
{} \& D_0 \& {}
\end{tikzcd}}
peut être caractérisé de diverses façons, p. ex

\diagram{\begin{tikzcd}[ampersand replacement=\&,column sep=.75em,row sep=.8em,execute at end picture={
\path (b.west) -- (c.west) coordinate[midway] (mid);
\draw ([xshift=-2pt]a.north west) .. controls +(-4pt,0pt) and +(0pt,3pt) .. ([xshift=-5pt]a.west) -- ([xshift=-5pt,yshift=4pt]mid) .. controls +(0pt,-3pt) and +(3pt,0pt) .. ([xshift=-9pt]mid) .. controls +(3pt,0pt) and +(0pt,3pt) .. ([xshift=-5pt,yshift=-4pt]mid) -- ([xshift=-5pt]d.west) .. controls +(0pt,-3pt) and +(-4pt,0pt) .. ([xshift=-2pt]d.south west);}]
a) \& |[alias=a]|\Lambda1, \& \Lambda2, \& \Lambda3,\arrow[d,phantom,"\unsure{\Updownarrow|\Downarrow}" description] \& C\Lambda\mathord{*}\arrow[d,Leftrightarrow]\\
b) \& |[alias=b]|\Lambda1, \& \Lambda2, \& D, \& C\mathord{*}\\
c) \& |[alias=c]|\Lambda1, \& C\Lambda\mathord{*}, \& D \& {}\\
d) \& |[alias=d]|\Lambda1,\unsure{''} \& \deleted{\unsure{C\mathord{*}}}\Lambda2, \& \deleted{\unsure{D_0}}C\mathord{*}, \& D_0
\end{tikzcd}}
\margin{les équivalences marquées \unsure{valent|valant} terme à terme, en présence de $\Lambda1,\ \Lambda2$ i.e. si $\Lambda$ est partie très dense de $M$}
\margin{\uline{NB} $\Lambda1$ est le seul des axiomes partiels qui figure dans \unsure{chacun} des quatre \unsure{syst.} équivalents a) b) c) d)}

On dira, quand ces conditions sont vérifiées qu'on a un atlas $\Lambda$-spacieux, ou $\Lambda$-polyédral. Les cas les plus intéressants pour moi sont

$1^{\circ})$ $\Lambda=L$ cas déjà traité --- mais cela suppose que l'atlas satisfait notamment la condition de divisibilité

$2^{\circ})$ $\Lambda=M$, auquel cas $\Lambda1$ \deleted{\unsure{est}} devient tautologique, $\Lambda2$ équivaut à $D_0$, $\Lambda3$ équivaut à $D$, $C\Lambda\mathord{*}$ équivaut à $C\mathord{*}$, \deleted{et la} Dire que $A$ est $M$-spacieux \inserted{ou $M$-polyédr\unsure{.}} signifie donc simplement qu'il satisfait \unsure{:|à} $D$, et à $C$ spac \unsure{resp} $C$ pol.

Dans le cas général, on peut dire que $A$ est (disons) $\Lambda$-spacieux, \unsure{dès|ssi|si} il est $M$-spacieux i.e. spacieux tout court, et si de plus $\Lambda$ est très dense ($\Lambda1$ et $\Lambda2$)
\scanpage{48}\mspage{4\deleted{\unsure{6}}7}

On voudrait donner un critère sur
\[
(M,\underset{M}{\leq},\underset{M}{\ll}\deleted{\unsure{,}},\Lambda\subset M,\underset{\Lambda}{\parallel})
\]
où $\underset{\Lambda}{\parallel}$ est une relation dans $\Lambda$, \inserted{partie fermée de $M$,} pour que ça provienne d'un atlas $\Lambda$-spacieux, resp. $\Lambda$-pol. On va, pour fixer les idées, prendre le cas $\Lambda$-spacieux.

\uline{$M\Lambda$}$0$, \uline{$M\Lambda$}$1$, \uline{$M\Lambda2$}, \uline{$M\Lambda3$} comme avant (p.42)

\uline{$M\Lambda4$} : La relation $\underset{\Lambda}{\parallel}$ (ou $\parallel$) sur $\Lambda$ est \deleted{réflexive} symétrique, antiréflexive $(X\parallel Y\Longrightarrow X\ne Y)$, et
\[
X\parallel Y,\quad X'\ll X,\quad Y'\ll Y\quad\Longrightarrow\quad X'\parallel Y'.
\]

\uline{$M\Lambda5$} Soient $X,Y\in M$ tels qu'il existe $Z\in M$ avec $X,Y\leq Z$, \deleted{$\unsure{L=\widetilde X\cap}\unsure{Y'|\widetilde Y}$} $X'\in\mathrm{omb}_{\Lambda}(X)$, $Y'\in\mathrm{omb}_{\Lambda}(Y)$,
\margin{comparer $ML\deleted{\unsure{4}}5$ p.34}
Posons
\[
\left\{\begin{aligned}
&L=\widetilde X\cap\widetilde Y\\
&X'_L=\widetilde{X'}\cap\mathrm{Omb}_{\deleted{\unsure{\Lambda}}}(L)\qquad Y'_L=\widetilde{Y'}\cap\mathrm{Omb}(L)\\
&\text{où }\mathrm{Omb}(L)\overset{\mathrm{def}}{=}\bigcup_{Z\in L}\mathrm{Omb}(Z)
\end{aligned}\right.
\]
% slanted
\margin{On en aura besoin sous forme plus forte, en supposant seulement $X\mathrel{\rlap{$>$}<} Y$ (au sens ci-dessous).}

Supposons $X'_L\parallel Y'_L$ au sens : $\forall\,X''\in X'_L$, $Y''\in Y'_L$, \inserted{on a} $X''\parallel Y''$. (\uline{NB} on aura $X''\leq X'$, $Y''\leq Y'$, donc $X'',Y''\in\Lambda$) Alors $X'\parallel Y'$.

% slanted
\margin{\uuline{Non}, cf. plus bas p.49}
Je dis \uline{que ça suffit} (c'est évidemment nécessaire). \deleted{Comme \unsure{tantôt}} On pose, pour $X,Y\in M$, si $L=\widetilde X\cap\widetilde Y$
\[
\begin{aligned}
X\underset{\deleted{\unsure{M}}}{\mathrel{\rlap{$>$}<}}Y&\Longleftrightarrow
\forall\,\substack{X'\in\mathrm{omb}_{\Lambda}(X)\\Y'\in\mathrm{omb}_{\Lambda}(Y)}
\quad(X'_L=\varnothing,Y'_L=\varnothing)\Longrightarrow X'\underset{\Lambda}{\parallel}Y'\\
X\underset{M}{\parallel}Y&\Longleftrightarrow
\forall\, X'\in\mathrm{omb}_{\Lambda}(X),Y'\in\mathrm{omb}_{\Lambda}(Y),\text{ on a } X'\underset{\Lambda}{\parallel}Y'.
\end{aligned}
\]
\scanpage{49}\mspage{48}

Ces relations sont symétriques, \unsure{.} $\mathrel{\rlap{$>$}<}$ est réflexive (tautologique), $\underset{M}{\parallel}$ antiréflexive (utiliser le fait que $\underset{\Lambda}{\parallel}$ l'est, et $M\Lambda2$, impliquant que $\mathrm{omb}_{\Lambda}(X)\ne\varnothing$). Je dis

\uline{Lemme 1} \unsure{Soient|Soit} $X,Y\in M$.
\[
\begin{aligned}
\text{a)}\quad X\underset{M}{\parallel}Y &\Longleftrightarrow X\mathrel{\rlap{$>$}<} Y,\quad\text{et}\quad \underbrace{\widetilde X\cap\widetilde Y}_{L}=\varnothing\\
\text{b)}\quad\text{Si } X,Y\in\Lambda,\quad&\text{alors}\quad X\underset{M}{\parallel}Y\Longleftrightarrow X\underset{\Lambda}{\parallel}Y
\end{aligned}
\]
(on omettra donc par la suite les indices $M,\Lambda$ dans $\parallel$)

Prouvons a). \deleted{\unsure{On a} $\Longrightarrow$ : $X\mathrel{\rlap{$>$}<} Y$ \unsure{trivial sur} définition, $\widetilde X\cap\widetilde Y=\varnothing$}

Prouvons $\Longrightarrow$ i.e. supposons $X\underset{M}{\parallel}Y$, prouvons :
\[
\begin{aligned}
1^{\circ})\quad &X\mathrel{\rlap{$>$}<} Y:\ \text{tautologique}\\
2^{\circ})\quad &\widetilde X\cap\widetilde Y=\varnothing,\quad\text{si on avait }Z\in\widetilde X\cap\widetilde Y,\ \text{soit}
\end{aligned}
\]
$Z'\in\mathrm{omb}_{\Lambda}(Z)$, alors $Z'\in\mathrm{omb}_{\Lambda}(X)\cap\mathrm{omb}_{\Lambda}(Y)$ \deleted{\unsure{donc par}.} $X\parallel Y$ on doit avoir $Z'\parallel Z'$, absurde (antiréflexivité).

Prouvons $\Longleftarrow$ i.e. supposons $X\mathrel{\rlap{$>$}<} Y$, $\widetilde X\cap\widetilde Y=\varnothing$. et prouvons $X\parallel Y$. \inserted{Mais comme \deleted{par déf.}} \deleted{Soit $X'\in\mathrm{omb}_{\Lambda}(X)$, $Y'\in\mathrm{omb}_{\Lambda}(Y)$.} \deleted{, de $X\mathrel{\rlap{$>$}<} Y$, \unsure{donc le}} \unsure{on|a} $L=\widetilde X\cap\widetilde Y=\varnothing$, \deleted{\unsure{alr|alors}} donc $\mathrm{omb}_{\Lambda}(L)=\varnothing$, \deleted{\unsure{elle se réduit à|il en résulte} : $\forall X'\in\mathrm{omb}_{\Lambda}(X)$} la définition de $X\mathrel{\rlap{$>$}<} Y$ montre que ça \unsure{signifie|signifi:} $X\parallel Y$.

\uline{Prouvons} b Résulte tautologiquement de $M\Lambda4$

\uline{Terminologie}\unsure{.} Soi\deleted{\unsure{t}}\unsure{ent|t} $P,Q$ deux parties de $M$. On pose
\[
\begin{aligned}
P\parallel Q&\overset{\mathrm{def}}{\Longleftrightarrow}\quad\forall X\unsure{\mathpunct{:}}\in P,\ \forall Y\in Q,\ \text{\inserted{on a}}\quad X\parallel Y\\
&\Longleftrightarrow\quad\forall X'\in\mathrm{omb}_{\Lambda}P,\ Y'\in\mathrm{omb}_{\Lambda}(Q),\ \text{\inserted{on a}}\quad X'\parallel Y'
\end{aligned}
\]
(où $\mathrm{omb}_{\Lambda}(P)=\bigcup_{X\in P}\mathrm{omb}_{\Lambda}(X)$, et \unsure{item|idem} pour $\mathrm{omb}_{\Lambda}(Q)$)
